% Rubric: Results.
% Report qualitative and quantitative findings accurately, with well-chosen
% figures and tables that have informative captions, labels and units.
% Save figures to figures/ and include them with \includegraphics, as in the
% commented example below.

\section{Results}
\label{sec:results}

\subsection{RQ1: targeted removal exposes corridor bottlenecks}
Figure~\ref{fig:robustness-cascade}a shows earlier fragmentation under targeted
station removal. The interpolated removal fraction at which the mean largest
component reaches half of the intact 86 stations is
\RandomCrossPercent\% for random removal, \DegreeCrossPercent\% for a fixed
degree ranking and \BetweenCrossPercent\% for a fixed betweenness ranking.
These estimates are distinct from the susceptibility peaks of the
second-largest component at 6\%, 2\% and 1\%, and depend on station-count
rounding and grid resolution; on this small tree, neither diagnostic
establishes a thermodynamic phase transition. Betweenness identifies
corridor-separating positions that degree alone cannot distinguish: many
internal stations have degree two but sit between different amounts of the
network.

\subsection{RQ2: tolerance depends on the load-update assumption}
For the maximum-load trigger, Oats Street, static capacity-weighted
redistribution at $\alpha=0.2$ closes all \MaxLoadFailed{} stations, including
the trigger, in \MaxLoadRounds{} secondary rounds. On the 0.025 grid, the first
sampled tolerance containing the closure is $\alpha=\CapacityAlphaStar$;
equal redistribution requires $\alpha=\EqualAlphaStar$
(Figure~\ref{fig:robustness-cascade}b). These are sampled thresholds for the
specified trigger and rule, not universal station-capacity requirements.
Dynamic unweighted betweenness produces no secondary failures on the frozen
tree (Section~\ref{sec:model}): the disconnected graph can still lose service.
The 2,000 random-trigger observations include zero secondary failures; cascade
rounds are not separate trials. The observed size
distribution is concentrated at a few finite outcomes
(Figure~\ref{fig:uncertainty}b).

\subsection{RQ3: service recovery and rail damage can move in opposite directions}
After a Bayswater closure at $\alpha=0.2$, the rail-only scenario serves
\BayRailServicePercent\% of the original 3,655 terminal pairs. Full candidate
bus standby serves \BayBusServicePercent\%, but closes two rail facilities
instead of one (Figure~\ref{fig:bus-mechanism}a). At the first load check,
Fremantle's load-to-capacity ratio rises from \FremantleRailRatio{} without
buses to \FremantleBusRatio{} with the full pool
(Figure~\ref{fig:bus-mechanism}b), above the fixed rail-only capacity, so it
fails next; bus paths still keep its terminal connected. At $\alpha=1.0$, the
full-pool Bayswater case serves all pairs with only the trigger closed, so
extra capacity can contain the rerouting overload that the restored paths
created.

The complete pool does not worsen served service in the stored scan;
a budgeted strategy can. After the three exclusive Airport Line facilities
close at $\alpha=0.2$, the existing-bus policy selects two links and additionally
closes West Leederville. Served terminal pairs fall from
\AirportRailServicePercent\% with no backup to \AirportBusServicePercent\%.
Endpoint-weighted served demand likewise falls from
\AirportRailDemandPercent\% to \AirportBusDemandPercent\%
(Table~\ref{tab:airport}). These weights sum scheduled stop counts at pair
endpoints, not observed OD demand. Travel-time penalties are conditional on
reachable pairs; comparing means alone could mistake lost long journeys for
faster service.

\subsection{RQ4: numerical uncertainty and model sensitivity}
The tested nested-prefix criteria recommend \RecommendedSeeds{} random trials:
the worst normalised mean change is \WorstMeanChange{} and the largest pointwise
95\% interval half width is \WorstHalfWidth{}, below the chosen 0.02 and 0.03
tolerances (Figure~\ref{fig:uncertainty}a). These criteria apply to the stored
conditions, not every transport model. Coarsening the capacity-rule alpha grid
from 0.025 to 0.1 shifts first sampled containment from
\CapacityAlphaStar{} to \CoarseAlphaStar{}. The demand/frequency-capacity
scenario gives \DemandAlphaStar{}; it changes the load proxy, capacity
reference and trigger, so it is a model comparison, not an isolated demand
effect. Frequency weighting alone leaves unique tree paths unchanged.

The bounded power-law test uses maximum-likelihood fitting, cutoff selection
and refitted simulation, not log-log regression
\citep{clauset2009power}. At static alpha 0.1--0.2 there are only two distinct
positive secondary sizes, making the tail ineligible under the three-size
requirement. At 0.3 the fitted hypothesis is rejected
($p=\TailPValue$, 500 refitted simulations), with its exponent at the lower
search boundary. Dynamic samples contain only zeros. The experiments therefore
do not support a power-law avalanche claim.

% \begin{figure}[t]
%   \centering
%   \includegraphics[width=0.8\linewidth]{figures/example.png}
%   \caption{State what is shown, under which conditions, and the units.}
%   \label{fig:example}
% \end{figure}

% Then refer to Figure~\ref{fig:example} in the discussion of the finding.
